\section*{Hoofdstuk \ref{ch:b3:developmental-genetics} --- Ontwikkelingsgenetica en de evolutie van de vorm}

\begin{solution}{exo:b3:developmental-genetics:1}
Maternaal effect: \emph{bicoid}, embryo's van mutante moeders missen
kop en borststuk. Gap: \emph{Krüppel}, een aaneengesloten blok
segmenten ontbreekt. Paarregel: \emph{fushi tarazu} (of
\emph{even-skipped}), elk tweede segment ontbreekt. Segmentpolariteit:
\emph{hedgehog} (of \emph{engrailed}), waarbij een deel van elk segment
is vervangen door een spiegelbeeld van de rest.
\end{solution}

\begin{solution}{exo:b3:developmental-genetics:2}
Een gen waarvan het product door de moeder in het ei wordt gelegd,
zodat de verschijningsvorm van het embryo haar genotype volgt. Haar
embryo's missen kop en borststuk en dragen aan beide uiteinden
achterste structuren; zij sterven. Het wildtype-allel van de vader zit
in de zygote, maar het mRNA van \emph{bicoid} wordt tijdens de
eicelvorming door de voedstercellen van de moeder gemaakt en vóór de
bevruchting aan de pool verankerd; de eigen kopie van de zygote zou uren
later worden afgeschreven, wanneer de tijd van de gradiënt voorbij is.
\end{solution}

\begin{solution}{exo:b3:developmental-genetics:3}
\omterm{def:b3:developmental-genetics:hox}{Colineariteit}: de volgorde van de \omterm{def:b3:developmental-genetics:hox}{Hox-genen} op het chromosoom komt
overeen met de volgorde van hun expressiegebieden langs het lichaam.
Achterste overwicht: waar er meerdere tot expressie komen, bepaalt het
achterste de identiteit. Een muis zonder \emph{Hoxc8} vertoont een
transformatie naar voren --- haar eerste lendenwervel neemt de
identiteit van een borstwervel aan en draagt een rib.
\emph{Antennapedia} in de kop maakt van antennen poten, de identiteit
van het tweede borstsegment.
\end{solution}

\begin{solution}{exo:b3:developmental-genetics:4}
Een getransplanteerde rugzijdige lip wekte een tweede volledige as op,
grotendeels uit cellen van de gastheer: het transplantaat gaf zijn
buren instructies --- een \omterm{def:b3:developmental-genetics:organiser}{organisator}. Zijn moleculen zijn
uitgescheiden tegenspelers van BMP (Chordin, Noggin, Follistatine).
Ectoderm wordt neuraal wanneer het BMP-signaal ontbreekt en opperhuid
wanneer BMP aanwezig is, dus is neuraal wat ectoderm doet wanneer het
niets te horen krijgt: de \omterm{def:b3:developmental-genetics:organiser}{organisator} werkt door een instructie het
zwijgen op te leggen.
\end{solution}

\begin{solution}{exo:b3:developmental-genetics:5}
$k = \ln 2/(35\times 60) = \qty{3.3e-4}{s^{-1}}$; $\lambda = \sqrt{4/
3.3\times 10^{-4}} = \qty{110}{\micro m}$; verhouding $\mathrm{e}^{250/110}
= \mathrm{e}^{2.27} = 9.7$.
\end{solution}

\begin{solution}{exo:b3:developmental-genetics:6}
Vier kopieën verdubbelen $c_{0}$: de grens verschuift met $\lambda\ln 2 =
\qty{69}{\micro m}$ naar achteren, tot \qty{309}{\micro m}; één kopie
halveert haar: \qty{69}{\micro m} naar voren, tot \qty{171}{\micro m}.
Elke verschuiving is \qty{14}{\%} van de eilengte --- meer dan wordt
waargenomen, wat zelf een aanwijzing is voor mechanismen die de
afhankelijkheid van de dosis dempen.
\end{solution}

\begin{solution}{exo:b3:developmental-genetics:7}
$\delta x = \lambda\,\delta c/c$: voor \qty{8}{\micro m} is $\delta c/c =
8/100 = \qty{8}{\%}$. Een Poisson-telling $N$ heeft relatieve fout
$1/\sqrt{N}$, dus moeten er $N = 1/0.08^{2} \approx 160$ moleculen
worden geteld --- enkele procenten van de moleculen die in een kern
aanwezig zijn, of één telling die over de tijd wordt opgeteld.
\end{solution}

\begin{solution}{exo:b3:developmental-genetics:8}
$k = \ln 2/35 = \qty{0.0198}{min^{-1}}$: $1 - \mathrm{e}^{-kt}$ geeft
$0.70$ na \qty{60}{min}, $0.83$ na \qty{90}{min}, $0.91$ na
\qty{120}{min} en $0.95$ na \qty{150}{min}. De gradiënt stijgt terwijl
hij wordt afgelezen nog met enkele procenten, wat een grens over het
afleesvenster met $\lambda\ln(1/0.9) \approx \qty{10}{\micro m}$ zou
verschuiven, ongeveer één kern: het embryo moet ofwel op een vast
tijdstip aflezen ofwel een aflezing gebruiken die ongevoelig is voor het
algemene niveau.
\end{solution}

\begin{solution}{exo:b3:developmental-genetics:9}
Een tweede ZPA geeft een spiegelbeeldige verdubbeling, 4-3-2-2-3-4: de
voorste cellen zien nu veel Shh en worden achterste vingers. Een
bolletje dat weinig Shh afgeeft geeft een gedeeltelijke verdubbeling
--- een extra vinger 2 vooraan, zeg 2-2-3-4 --- omdat de voorste cellen
alleen de lage concentratie zien die vinger 2 vastlegt.
\end{solution}

\begin{solution}{exo:b3:developmental-genetics:10}
Bijzondere oplossing $s/k$; homogene oplossingen $\cosh(x/\lambda)$,
$\sinh(x/\lambda)$; geen stroom bij $0$ doodt de $\sinh$; $c(L) = 0$
legt de amplitude vast: $c(x) = (s/k)\bigl[1 - \cosh(x/\lambda)/\cosh(L/\lambda)\bigr]$.
Het profiel is een plateau nabij $x = 0$ dat in een grenslaag ter
breedte $\lambda$ bij de put tot nul daalt --- geen e-macht vanuit een
punt, maar een vlak veld met een rand; de hoogste concentratie ligt het
verst van de put. Voor $L \ll \lambda$ geldt $\cosh u \approx 1 +
u^{2}/2$ en $c \approx s(L^{2} - x^{2})/(2D)$: een parabool die niet van
$k$ afhangt, omdat de moleculen de put bereiken voordat zij kunnen
vervallen.
\end{solution}

\begin{solution}{exo:b3:developmental-genetics:11}
$x_{1} = \lambda\ln(1/0.3) = 1.20\lambda$ en $x_{2} = \lambda\ln(1/0.08)
= 2.53\lambda$: het eerste gebied is $1.20\lambda$ breed, het tweede
$1.32\lambda$, beide onafhankelijk van $c_{0}$, dat ze samen
verschuift. Het derde gebied, $L - 2.53\lambda$, vangt alle variatie in
de eilengte op: in een ei van \qty{450}{\micro m} met
$\lambda = \qty{100}{\micro m}$ is het achterlijf \qty{197}{\micro m},
in een ei van \qty{550}{\micro m} \qty{297}{\micro m} --- het patroon
schaalt niet mee. Mechanismen: een tweede gradiënt vanaf de achterste
pool (Nanos, Caudal, of het terminale systeem) die samen met Bicoid
wordt gelezen, zodat de posities door de verhouding worden vastgelegd;
of een ``uitrekker'' die wordt gemaakt waar het morfogeen laag is en
$\lambda$ verlengt tot de gradiënt het ei vult.
\end{solution}

\begin{solution}{exo:b3:developmental-genetics:12}
Een verandering in de codering verandert het eiwit overal waar het
werkt: \emph{Pitx1} bouwt ook de \omterm{def:b3:endocrinology:axes}{hypofyse} en de kaak, \emph{\omterm{def:b3:developmental-genetics:organiser}{sonic
hedgehog}} patroneert ook de neurale buis, de darm en de tanden, dus is
een kapot eiwit dodelijk of pleiotroop verlammend. Een versterker
stuurt de expressie in één weefsel op één tijdstip; hem verwijderen
haalt de structuur weg die dat weefsel zou hebben gebouwd en laat elk
ander gebruik van het eiwit ongemoeid. Versterkers zijn modulair en
vaak ten dele overtollig, zodat tussenstappen leefbaar zijn, en
dezelfde weg --- een verloren schakelaar --- is in vele
stekelbaarsmeren onafhankelijk ingeslagen.
\end{solution}

\begin{solution}{pb:b3:developmental-genetics:1}
\textbf{1.} $k = \ln 2/(35\times 60) = \qty{3.3e-4}{s^{-1}}$; $1/k =
\qty{3030}{s} = \qty{50}{min}$.
\textbf{2.} $\lambda = \sqrt{4/3.3\times 10^{-4}} = \qty{110}{\micro m}$;
het ei is $4.5\lambda$.
\textbf{3.} $\mathrm{e}^{4.5} \approx 90$ van pool tot pool;
$\mathrm{e}^{8.5/110} = 1.08$, een verandering van \qty{8}{\%} per kern.
\textbf{4.} $J = 50\times\sqrt{4\times 3.3\times 10^{-4}} = 50\times
0.036 = \qty{1.8}{nM\,\micro m\,s^{-1}}$. Eén nM is $0.6$ moleculen per
\unit{\micro m^3}, dus ongeveer $1.1$ moleculen per \unit{\micro m^2}
per seconde.
\textbf{5.} Kernvolume $\tfrac{4}{3}\pi\times 27 = \qty{113}{\micro
m^3}$. Bij \qty{50}{nM} zijn dat $30$ moleculen per \unit{\micro m^3}:
\num{3400} moleculen. Halverwege het ei is $c = 50\,\mathrm{e}^{-2.27} =
\qty{5.2}{nM}$: ongeveer \num{350} moleculen.
\textbf{6.} $kt = 1.19$ na \qty{60}{min}: $0.70$; $kt = 2.38$ na
\qty{120}{min}: $0.91$. De gradiënt zit bij het aflezen op
\qty{10}{\%} van de stationaire waarde --- nog altijd stijgend met een
bedrag dat een grens met $\lambda\ln(1/0.91) \approx \qty{10}{\micro m}$
zou verplaatsen, één kern.
\textbf{7.} $x_{1} = 110\ln(1/0.3) = \qty{132}{\micro m}$
(\qty{26}{\%}); $x_{2} = 110\ln(1/0.08) = \qty{278}{\micro m}$
(\qty{56}{\%}).
\textbf{8.} Vier kopieën: $c_{0} = \qty{100}{nM}$; beide grenzen
verschuiven met $\lambda\ln 2 = \qty{76}{\micro m}$ naar achteren, tot
$208$ en \qty{354}{\micro m}. Zes kopieën: $\lambda\ln 3 = \qty{121}{\micro m}$,
tot $253$ en \qty{399}{\micro m}.
\textbf{9.} Eén kopie: \qty{76}{\micro m} naar voren, tot $56$ en
\qty{202}{\micro m}. Het kopgebied kromp van $132$ tot
\qty{56}{\micro m}; het borststuk behield zijn breedte
(\qty{146}{\micro m}) en schoof naar voren; het achterlijf groeide van
$222$ tot \qty{298}{\micro m}.
\textbf{10.} $\delta x = 110\times 0.1 = \qty{11}{\micro m}$, $1.3$
kernafstanden.
\textbf{11.} Eén kern: $\delta c/c = 8.5/110 = \qty{7.7}{\%}$. Vanaf
\qty{10}{\%} is $n = (10/7.7)^{2} \approx 1.7$: twee onafhankelijke
aflezingen (of het gemiddelde van twee naburige kernen) volstaan.
\textbf{12.} $1/\sqrt{350} = \qty{5.3}{\%}$: één ogenblikkelijke
telling is al nauwkeuriger dan de aangenomen \qty{10}{\%}, dus komt de
waargenomen ruis niet alleen van het tellen maar ook van de
bindingskinetiek en de afleesmachinerie --- en middelen over de tijd
verbetert haar nog.
\textbf{13.} $278/450 = \qty{62}{\%}$ en $278/550 = \qty{51}{\%}$ van de
eilengte: een verschil van \qty{11}{\%}, zes kernen. Het patroon
schaalt niet mee bij een vaste $\lambda$.
\textbf{14.} Als $D \propto L^{2}$ bij een vaste $k$, dan is
$\lambda \propto L$ en is $x_{T}/L = (\lambda/L)\ln(c_{0}/T)$ in elk ei
hetzelfde: volmaakte schaling. Maar $D$ is een eigenschap van het
molecuul en het cytoplasma, niet van de grootte van het ei, dus is dit
zoals gesteld onaannemelijk; de schaling in echte embryo's is
gedeeltelijk en komt van andere aflezingen.
\textbf{15.} $c^{5}/(c^{5} + T^{5}) = 0.1$ bij $c/T = 9^{-1/5} = 0.64$
en $0.9$ bij $c/T = 9^{1/5} = 1.55$: een factor $2.4$ in concentratie,
$\Delta x = 110\ln 2.4 = \qty{97}{\micro m}$ --- elf kernen, veel
breder dan de waargenomen grens van twee of drie; het aflezen van de
gradiënt alleen is niet scherp genoeg.
\textbf{16.} Steilheid zet een gegeven concentratie beslissender om in
aan of uit, zodat het overgangsgebied smaller is; maar een fout van
\qty{10}{\%} in de concentratie verplaatst het punt waar de drempel
wordt overschreden nog altijd met $\lambda\,\delta c/c$: het antwoord
kan een volmaakte stap zijn en die stap toch op de verkeerde plek
landen.
\textbf{17.} Stonden beide aan, dan zou elk het andere onderdrukken,
zodat er ten minste één zou worden uitgezet; bij coöperatieve
onderdrukking zijn de enige stabiele toestanden één aan en één uit, en
is een grens ertussen een schakelaar en geen helling.
\textbf{18.} $\mathrm{d}\ln\lambda = \tfrac{1}{2}(0.02 - 0.06) = -0.02$
per graad: $\lambda$ daalt \qty{2}{\%} per graad, \qty{10}{\%} bij
\qty{5}{\celsius}, tot \qty{99}{\micro m}; de grens van
\emph{hunchback} verschuift van $278$ naar \qty{250}{\micro m}, drie
kernen naar voren. Vliegen ontwikkelen zich tussen \qty{18}{} en
\qty{29}{\celsius} met normale verhoudingen, dus compenseert er iets
--- een bron of een aflezing die met de temperatuur de andere kant op
schuift.
\textbf{19.} $2^{8} = 256$ codes; er worden er veertien gebruikt. Bij
achterste overwicht telt alleen het achterste gen dat aan staat, dus
zijn er hoogstens negen verschillende identiteiten (acht genen, of
geen): de code is een rangorde, geen binair woord.
\textbf{20.} Het derde borstsegment wordt een tweede: \emph{Ubx}
onderdrukte het vleugelprogramma dat een borstsegment standaard
uitvoert, en de repressor weghalen maakt de voorouderlijke toestand
vrij --- vier vleugels, zoals bij de viervleugelige voorouders van de
vlieg.
\textbf{21.} Het halsgebied van de gans loopt tot wervel 22: een lange
nek van vele wervels, tegen zeven bij de muis. Het gen is hetzelfde; de
voorste grens van zijn expressie is naar achteren verschoven --- een
verandering in de regeling van waar een Hox-gen wordt aangezet.
\textbf{22.} Elke rand is nu een bron; Shh is aan beide randen het
hoogst (vinger 4), daalt naar het midden (3, dan 2) en bereikt nooit
het lage niveau dat vinger 1 vastlegt, zodat het midden twee vingers 2
rug aan rug bevat: 4-3-2-2-3-4.
\textbf{23.} \emph{Pax6} is een geconserveerde meesterregelaar die
welk oogprogramma de gastheer ook draagt in gang kan zetten: de vlieg
bouwde een vliegenoog, dus draagt het muizengen geen informatie over
een muizenoog, alleen het bevel ``bouw hier een oog''. Het toont de
\omterm{def:b3:developmental-genetics:evodevo}{diepe homologie} van het gennetwerk, niet de homologie van de organen:
het samengestelde oog en het cameraoog zijn afzonderlijk ontstaan uit
een voorouder met een plek fotoreceptoren onder \emph{Pax6}.
\textbf{24.} Vijf vingers over \qty{2}{mm} is een golflengte van
\qty{0.4}{mm}; zes vergt \qty{0.33}{mm}, een verkorting met een zesde.
Omdat $\lambda \propto \sqrt{D_{\text{inh}}/k_{\text{inh}}}$, verhoog
het verval van de remmer met \qty{44}{\%} of verlaag zijn diffusie met
\qty{31}{\%}; in de handplaat lukt dat door de dosis van de achterste
\omterm{def:b3:developmental-genetics:hox}{Hox-genen} te verlagen, en de mutante muizen hebben zes, zeven of meer
dunne vingers.
\textbf{25.} $\lambda \approx \qty{110}{\micro m}$; een verdubbeling
van de bron verschuift elke grens met \qty{76}{\micro m}; een ruis van
\qty{10}{\%} in de concentratie is \qty{11}{\micro m} positiefout,
ongeveer één kern.
\end{solution}
